\chapter{Recommended Priorities}
Based on the identified security controls stated in the chapter before, I have selected the following three controls that should receive the highest priority if Equifax were to implement them, starting with the most important one:
\begin{enumerate}
\item \textbf{Continuous Vulnerability Management:} This control addresses the risk  of known vulnerabilities remaining unpatched and being exploited by attackers. In the incident, the initial point of compromise was that the patch for the CVE-2017-5638 vulnerability was not applied to ACIS, allowing the attackers to gain access to Equifax's network. The assets this control protects are applications such as ACIS, servers and systems connected to them. It could have prevented or at least limited the unauthorized access to Equifax's internal systems and the subsequent exfiltration of consumer data. Therefore, it receives the highest priority because preventing the initial compromise can stop the following stages of the incident altogether. While other controls help to limit the impact of the incident, an effective vulnerability management can help to eliminate threats before they can be exploited.
\item \textbf{Network Infrastructure Management:} This control addresses the risk of attackers moving through a network after they have compromised one system. In the incident, the attackers were able to move through Equifax's network and access their database systems after compromising ACIS even though both systems were unrelated to each other. The assets this control protects are internal networks and all the applications and systems inside of these networks and can prevent or limit an attacker's ability to move through an organization's network and thus, reduce the number of systems they can reach. This security control receives the second highest priority because restricting access between systems can contain and severely limit the scope of a breach when preventive measures fail.
\item \textbf{Data Protection:} Finally, this control addresses the risk of data being accessed, manipulated or even exfiltrated by malicious actors. The assets protected  is sensitive consumer data and the control helps to prevent or limit the unauthorized access or theft of this data. Data protection is given a high priority as it represents the target of these breaches itself and measures such as encryption or access restriction can reduce the consequences of such an incident. However, it is prioritized below the other two controls as data protection itself does not prevent or contain attacks itself.
\end{enumerate}
To summarize, these controls provide complementary layers of protection. Vulnerability management reduces the likelihood of an compromise, network infrastructure management contains the breach after the initial compromise while data protection reduces the potential impact of an incident if sensitive data is accessed. 
\newline
Although these security control address the likelihood and potential consequences of an incident, it has to be noted that this incident in particular has been the result of a combination of technical and organizational failure. 
\newline
Specifically, prior to the incident, the security certificate for ACIS has been expired for 19 months, up to 300 were expired with 79 expired certificates belonging to business critical domains. Additionally, the security department that develops security policies and the IT department that implements these policies were divided, causing ineffective collaboration and coordination.
\newline
In that sense, although implementing these security controls would have helped Equifax, I believe that the problems lie in a weak cybersecurity culture within the organization and therefore, strengthening this culture should have the highest priority before implementing security controls.
\newline
Equifax should focus on developing an actual security culture rather than to implementing security controls. The people in the organization should understand, why certain measures and controls exists while management should ensure that security is integrated in every step of their operations as even the best technical measures and tools are as effective as the people implementing them, how they act, behave and perceive risks.